\section{Recommendations}\label{sec:recs}
\paragraph{Minimum detectable effect.}
For a paired comparison of two models on one test set of $n$ molecules, with per-molecule differences of squared errors $d_i=e_{A,i}^2-e_{B,i}^2$, the smallest MSE difference detectable at level $\alpha$ with power $1-\beta$ is approximately $(z_{1-\alpha/2}+z_{1-\beta})\,\mathrm{sd}(d)/\sqrt{n}$; dividing by $2\,\mathrm{RMSE}_B$ converts it to RMSE units (delta method), and dividing by $\mathrm{RMSE}_{\text{mean}}-\mathrm{RMSE}_{\text{floor}}$ to units of $f$. On our test sets (200--1{,}900 molecules) the median MDE ($\alpha=0.05$, power 0.8) is about 0.04 in $f$ for ECFP+RF versus ECFP+GBM, whose errors are strongly correlated, and about 0.08 for frozen embeddings versus ECFP+GBM (\texttt{results/full\_clean/mde.csv}). Differences smaller than this on a single dataset should not be reported as improvements; detecting them requires pooling over tasks with a paired test, as we do here. The approximation assumes independent test molecules. Molecules from the same document have strongly correlated errors --- Section~\ref{sec:support} measures how much their error depends on shared document membership --- so these MDEs are optimistic, and more so on random splits, where documents are most heavily represented on both sides. It has not been validated by simulation.

\paragraph{Report document support, and split on it.}
Section~\ref{sec:support} shows that the number of training molecules drawn from a test molecule's own source document determines, more than the split type, how well a model scores: with none, test RMSE is about 1.05 pChEMBL units; with a hundred or more, about 0.51, below the estimated noise floor. Two consequences follow for anyone building or reading a potency benchmark.

First, a benchmark should report the distribution of document support over its test set, not only the split type. Two papers that both say ``scaffold split'' can differ by a factor of two in error for this reason alone, and the difference will be read as model quality. The distribution costs nothing to compute from the same table that produced the split.

Second, a split grouped by source document should be reported alongside the usual three. It is the only one of the four in which a test molecule's label comes from a document the model has never seen, which is the situation a prospective user faces. It is also cheap: it needs the document identifier that ChEMBL already supplies, and no chemistry. We do not propose it as a replacement for scaffold or temporal splits, which answer different questions, but a random split without it now tells us little, because the quantity it measures is dominated by how much of each document the model was given.

A corollary for the floor: a score at or past the noise ceiling on a random or scaffold split should be read as evidence of document support, not as a model approaching the limit of the data. The ceiling is a useful detector precisely because crossing it is diagnostic.

\paragraph{Reporting protocol.}
For each task, report: the estimated between-document standard deviation with confidence intervals from a bootstrap over document pairs, the number of replicated molecules and of distinct document pairs, the trimming rule, and how duplicated values and stereoisomers were handled; $f$ alongside RMSE and $R^2$; the ChEMBL release and retrieval date; split type and seeds; paired differences against a tuned ECFP baseline with a test over tasks; and the decontamination performed (exact, connectivity and near-analogue matches against public benchmarks; status of pretraining-corpus overlap).

\paragraph{Checklist.}
\begin{itemize}[nosep]
\item Are replicates from independent documents, and how many molecules have them?
\item Is the floor computed for the test molecules, with a confidence interval over document pairs and sensitivity to trimming and duplicated values?
\item Is the ECFP baseline tuned with the same budget as the proposed model?
\item Are random, scaffold, temporal \emph{and document-grouped} splits reported, with at least five seeds?
\item Is the distribution of document support over the test set reported for each split?
\item Is every $f>1$ investigated rather than reported as success, starting with document support (Section~\ref{sec:support}), then benchmark leakage and the floor estimate?
\item Is overlap of the test set with pretraining data checked, or stated as unknown?
\item Is the claimed improvement larger than the minimum detectable effect, or tested across tasks?
\end{itemize}
